%% =====================================================================
%%  TEMPLATE DE ARTIGO — ATIVIDADE DE REDES NEURAIS ARTIFICIAIS
%%  Modelo IEEE (classe IEEEtran) — usar no Overleaf
%% ---------------------------------------------------------------------
%%  INSTRUÇÕES RÁPIDAS:
%%  1. No Overleaf, clique em "New Project > Upload Project" e envie
%%     este arquivo (ou cole o conteúdo em um projeto novo).
%%  2. Substitua os textos de exemplo pelos seus resultados.
%%  3. Gere as figuras (matriz de confusão, curva ROC) em matplotlib e
%%     faça upload dos arquivos .png/.pdf para o Overleaf.
%%  4. Limite: 3 páginas no máximo (incluindo figuras e referências).
%% =====================================================================
\documentclass[conference]{IEEEtran}
\IEEEoverridecommandlockouts

% ---------- Pacotes (não é necessário adicionar mais) ----------
\usepackage{cite}
\usepackage{amsmath,amssymb,amsfonts}
\usepackage{algorithmic}
\usepackage{graphicx}
\usepackage{textcomp}
\usepackage{xcolor}
\usepackage{booktabs}
\usepackage{url}

% ---------- Metadados do PDF ----------
\usepackage[hidelinks]{hyperref}
% Permite que URLs longas quebrem de linha (evita linha ultrapassando a coluna)
\def\UrlBreaks{\do/\do.\do-\do:}
\hypersetup{
    pdftitle={Template de Artigo - Redes Neurais Artificiais},
    pdfauthor={Nome do Aluno 1, Nome do Aluno 2, Nome do Aluno 3}
}

\begin{document}

% ---------- Título e autores (EDITE AQUI) ----------
\title{Exemplo de Título: Comparação de Arquiteturas de MLP para
Classificação do Dataset Breast Cancer Wisconsin\thanks{Trabalho
apresentado à disciplina de Redes Neurais Artificiais.}}
% Dica: o título deve refletir a pergunta experimental do seu artigo.

\author{\IEEEauthorblockN{Nome do Aluno 1}
\IEEEauthorblockA{Departamento de Computação\\
Universidade Exemplo\\
Email: aluno1@exemplo.edu}
\and
\IEEEauthorblockN{Nome do Aluno 2}
\IEEEauthorblockA{Departamento de Computação\\
Universidade Exemplo\\
Email: aluno2@exemplo.edu}
\and
\IEEEauthorblockN{Nome do Aluno 3}
\IEEEauthorblockA{Departamento de Computação\\
Universidade Exemplo\\
Email: aluno3@exemplo.edu}}

\maketitle

% ---------- Resumo ----------
\begin{abstract}
% Resumo de 4 a 6 linhas: contexto, o que foi testado, método, principais
% resultados (com números) e conclusão. Exemplo:
Este artigo compara três arquiteturas de perceptron multicamadas (MLP)
na tarefa de classificação do dataset Breast Cancer Wisconsin. Foram
avaliadas redes rasa (16), média (32--16) e profunda (64--32--16),
mantendo fixos o otimizador Adam, a taxa de aprendizado e o número de
épocas. Os resultados mostram que a arquitetura média obteve o melhor
equilíbrio entre desempenho e generalização, com acurácia de 97,4\% e
F1-score de 0,974, enquanto a rede profunda apresentou sinais de
overfitting.
\end{abstract}

\begin{IEEEkeywords}
Redes neurais artificiais, perceptron multicamadas, classificação,
métricas de avaliação, matriz de confusão, curva ROC.
\end{IEEEkeywords}

% =====================================================================
\section{Introdução}
% Contextualize o problema (2 parágrafos) e, no final, faça a pergunta
% experimental. Exemplo:
A escolha da arquitetura de uma rede neural -- número de camadas e de
neurônios -- influencia diretamente sua capacidade de generalização.
Redes muito simples podem subajustar (underfitting), enquanto redes
muito complexas podem memorizar o conjunto de treino e falhar no teste
(overfitting)~\cite{goodfellow2016deep}.

Neste trabalho, investiga-se a seguinte pergunta: \emph{como a
profundidade e a largura de uma MLP afetam o desempenho de
classificação no dataset Breast Cancer Wisconsin?} A resposta é
verificada experimentalmente por meio das métricas estudadas em sala:
acurácia, precisão, recall, F1-score, curva ROC e matriz de confusão.

% =====================================================================
\section{Metodologia}

\subsection{Dataset}
% Descreva: origem com citação, número de amostras, atributos, classes e
% qualquer pré-processamento (normalização, dados faltantes, etc.).
Utilizou-se o dataset Breast Cancer Wisconsin (Diagnostic), disponível
no repositório UCI Machine Learning~\cite{kelly2023uci}. O dataset
possui 569 amostras, 30 atributos numéricos e 2 classes (maligno e
benigno). Os atributos foram normalizados para média 0 e desvio padrão
1. A divisão entre treino (80\%) e teste (20\%) foi estratificada por
classe, com semente aleatória fixada em 42 para reprodutibilidade.

\subsection{Modelo e protocolo experimental}
% Descreva arquitetura(s), ativação, perda, otimizador, épocas e protocolo.
Foram treinadas três MLPs com função de ativação ReLU nas camadas
ocultas e sigmoide na camada de saída, função de perda
binary cross-entropy e otimizador Adam~\cite{kingma2015adam} com taxa de
aprendizado $10^{-3}$, por 100 épocas. As arquiteturas avaliadas foram:
MLP rasa (16), MLP média (32--16) e MLP profunda (64--32--16). Todas as
redes foram implementadas com a biblioteca scikit-learn~\cite{pedregosa2011scikit}.

% =====================================================================
\section{Resultados e Discussão}
% Apresente as métricas em tabela, a matriz de confusão e a curva ROC.
% INTERPRETE os números: o que a matriz de confusão revela? A ROC?

A Tabela~\ref{tab:metricas} resume o desempenho das três arquiteturas no
conjunto de teste. A MLP média obteve a melhor acurácia (97,4\%) e
F1-score (0,974). A rede profunda teve acurácia de treino de 100\%,
mas desempenho inferior no teste, indicando overfitting.

\begin{table}[htbp]
    \caption{Desempenho das arquiteturas no conjunto de teste.}
    \label{tab:metricas}
    \centering
    \setlength{\tabcolsep}{4pt}
    \footnotesize
    \begin{tabular}{lcccc}
        \toprule
        \textbf{Arquitetura} & \textbf{Acurácia} & \textbf{Precisão} &
        \textbf{Recall} & \textbf{F1} \\
        \midrule
        Rasa (16)             & 0,965 & 0,960 & 0,958 & 0,959 \\
        Média (32--16)        & \textbf{0,974} & \textbf{0,976} &
        \textbf{0,972} & \textbf{0,974} \\
        Profunda (64--32--16) & 0,956 & 0,951 & 0,949 & 0,950 \\
        \bottomrule
    \end{tabular}
\end{table}

% Dica: substitua os valores acima pelos seus resultados reais.

A Figura~\ref{fig:matriz} apresenta a matriz de confusão da melhor
arquitetura. Observa-se apenas 1 falso negativo, o que é crítico em
aplicações médicas, pois representa um tumor maligno classificado como
benigno.

\begin{figure}[htbp]
    \centerline{\includegraphics[width=\columnwidth]
    {matriz_confusao.png}}  % <-- faça upload da sua figura no Overleaf
    \caption{Matriz de confusão da MLP média no conjunto de teste.}
    \label{fig:matriz}
\end{figure}

A curva ROC (Figura~\ref{fig:roc}) confirma o excelente poder de
discriminação do modelo, com área sob a curva (AUC) de 0,997.

\begin{figure}[htbp]
    \centerline{\includegraphics[width=\columnwidth]
    {curva_roc.png}}  % <-- faça upload da sua figura no Overleaf
    \caption{Curva ROC da MLP média no conjunto de teste.}
    \label{fig:roc}
\end{figure}

% =====================================================================
\section{Conclusão}
% Responda à pergunta da introdução, cite limitações e resultados.
Os resultados indicam que a arquitetura de profundidade intermediária
(32--16) ofereceu o melhor compromisso entre capacidade e
generalização, respondendo à pergunta proposta. Como limitação,
destaca-se o uso de uma única divisão treino/teste; trabalhos futuros
poderiam empregar validação cruzada para aumentar a robustez
estatística das comparações.

% =====================================================================
\begin{thebibliography}{00}
% Substitua/adicione suas referências (sempre cite o dataset e as
% bibliotecas utilizadas). Modelo IEEE: Autor, "Título", fonte, ano.
\bibitem{kelly2023uci}
M.~Kelly, R.~Longjohn, and K.~Nottingham, ``The UCI Machine Learning
Repository,'' University of California, Irvine, 2023. [Online].
Available: \href{https://archive.ics.uci.edu}{archive.ics.uci.edu}
\bibitem{goodfellow2016deep}
I.~Goodfellow, Y.~Bengio, and A.~Courville, \emph{Deep Learning}.
Cambridge, MA, USA: MIT Press, 2016.
\bibitem{kingma2015adam}
D.~P. Kingma and J.~Ba, ``Adam: A method for stochastic
optimization,'' in \emph{Proc. 3rd Int. Conf. Learning
Representations (ICLR)}, 2015.
\bibitem{pedregosa2011scikit}
F.~Pedregosa et al., ``Scikit-learn: Machine learning in Python,''
\emph{Journal of Machine Learning Research}, vol. 12, pp. 2825--2830,
2011.
\end{thebibliography}

\end{document}
